\documentclass[a4paper,12pt]{article}
\usepackage[utf8]{inputenc}
\usepackage[russian]{babel}
\usepackage{amsmath, amssymb, amsthm}
\usepackage{geometry}
\geometry{left=2cm,right=2cm,top=2cm,bottom=2cm}
\usepackage{enumitem}

\title{}
\author{}
\date{}

\begin{document}
	
	\maketitle
	
	\bigskip

	
	\section*{Основная часть}
	
	\subsection*{Задача 1 }
	
	
	\medskip
	
	Выпишем лагранжиан, вводя двойственные переменные:
	\begin{itemize}[leftmargin=1cm]
	 \item $u$ для неравенства $4x_1+5x_2+6x_3\le 7$, с условием $u\ge0$;
	 
	 \item $v$ для равенства $8x_1+9x_2+10x_3=11$, переменная $v\in\mathbb{R}$.
	\end{itemize}
	
	Лагранжиан имеет вид:
	\[
	\begin{array}{rl}
		L(x,u,v) &= x_1+2x_2+3x_3 + u\Bigl(7 - (4x_1+5x_2+6x_3)\Bigr) + v\Bigl(11 - (8x_1+9x_2+10x_3)\Bigr)\\[1mm]
		&= 7u + 11v + x_1\Bigl(1-4u-8v\Bigr) + x_2\Bigl(2-5u-9v\Bigr) + x_3\Bigl(3-6u-10v\Bigr).
	\end{array}
	\]
	Чтобы функция $L$ была ограничена сверху (причём в задаче максимизации нам нужен супремум по $x$ конечен), необходимо выполнить следующие условия:
	\begin{enumerate}
		\item Для переменной $x_1$, которая удовлетворяет условию $x_1\ge 0$, чтобы 
		\[
		\sup_{x_1\ge 0}\; x_1\Bigl(1-4u-8v\Bigr) < +\infty,
		\]
		требуется
		\[
		1-4u-8v\le 0.
		\]
		\item Для переменных $x_2$ и $x_3$, которые являются неограниченными (свободными), оптимальность требует равенств:
		\[
		2-5u-9v = 0,\quad 3-6u-10v = 0.
		\]
	\end{enumerate}
	Таким образом, двойственная задача имеет вид:
	\[
	\begin{array}{rl}
		\text{minimize} & 7u+11v\\[1mm]
		\text{s.t.} & 1-4u-8v \le 0,\\[1mm]
		& 5u+9v = 2,\\[1mm]
		& 6u+10v = 3,\\[1mm]
		& u\ge 0,\quad v\in \mathbb{R}.
	\end{array}
	\]
	\vspace{2mm}
	
	\textbf{Решение системы:} Из равенств:
	\[
	\begin{cases}
		5u+9v=2,\\[1mm]
		6u+10v=3,
	\end{cases}
	\]
	вычтем, домножив первое равенство на 6 и второе на 5:
	\[
	\begin{array}{rcl}
		30u+54v&=&12,\\[1mm]
		30u+50v&=&15.
	\end{array}
	\]
	Вычитая второе у первого, получаем:
	\[
	4v = -3 \quad\Longrightarrow\quad v = -\frac{3}{4}.
	\]
	Подставляя $v$ в первое равенство:
	\[
	5u+9\Bigl(-\frac{3}{4}\Bigr)=2\quad\Longrightarrow\quad 5u - \frac{27}{4} =2,\quad 5u = \frac{35}{4},\quad u = \frac{7}{4}.
	\]
	Проверка неравенства:
	\[
	1-4\Bigl(\frac{7}{4}\Bigr)-8\Bigl(-\frac{3}{4}\Bigr)= 1-7+6 =0.
	\]
	Таким образом, оптимальное решение двойственной задачи: $u=\frac{7}{4}$, $v=-\frac{3}{4}$, а оптимальное значение равно:
	\[
	7u+11v=7\cdot\frac{7}{4}+11\Bigl(-\frac{3}{4}\Bigr)=\frac{49-33}{4}=\frac{16}{4}=4.
	\]
	
	\bigskip
	
	\subsection*{Задача 2}
	
	
	\medskip 
	\begin{enumerate}[label=\arabic*)]
		\item Первое ограничение: 
		\[
		-4x_1+4x_3\le 1,\quad \Rightarrow\quad \text{двойственный множитель } y_1\ge 0.
		\]
		\item Второе ограничение:
		\[
		5x_1+x_2+7x_3=0,\quad \Rightarrow\quad y_2\in\mathbb{R}.
		\]
		\item Третье ограничение задано как
		\[
		-8x_1-7x_2+3x_3 \ge 2.
		\]
		Домножим на $-1$:
		\[
		8x_1+7x_2-3x_3 \le -2,\quad \Rightarrow\quad \text{двойственный множитель } y_3\ge 0.
		\]
		\item Четвёртое ограничение:
		\[
		6x_1+3x_2+8x_3 \le 3,\quad \Rightarrow\quad y_4\ge 0.
		\]
	\end{enumerate}
	
	Запишем двойственную задачу, пользуясь стандартной схемой для задачи минимизации:
	\[
	\begin{array}{rl}
		\text{maximize} & y_1\cdot 1 + y_2\cdot 0 + y_3\cdot (-2) + y_4\cdot 3 = y_1 -2y_3+3y_4\\[1mm]
		\text{s.t.} & \text{для каждой координаты } x_j\text{ должно выполниться } (A^T y)_j \,\Diamond\, c_j,
	\end{array}
	\]
	где знак неравенства ($\Diamond$) выбирается исходя из ограничений на переменные:
	\begin{itemize}[leftmargin=1cm]
		\item Если $x_j \ge 0$, то $(A^T y)_j \le c_j$;
		\item Если $x_j$ свободна, то $(A^T y)_j = c_j$.
	\end{itemize}
	
	Определим коэффициенты для каждой переменной:
	\begin{itemize}[leftmargin=1cm]
		\item Для $x_1$: коэффициенты в ограничениях:
		\[
		\begin{array}{ll}
			\text{Ограничение 1:} & -4,\\[1mm]
			\text{Ограничение 2:} & 5,\\[1mm]
			\text{Ограничение 3 (после домножения):} & 8,\\[1mm]
			\text{Ограничение 4:} & 6.
		\end{array}
		\]
		Тогда требуемое соотношение:
		\[
		-4y_1 + 5y_2 + 8y_3 + 6y_4 \ge c_1,
		\]
		где $c_1=-9$. Однако поскольку $x_1\ge 0$ (по условию) условие должно иметь вид:
		\[
		-4y_1 + 5y_2 + 8y_3 + 6y_4 \le -9.
		\]
		Это означает, что при переходе к двойственной задаче знак неравенства обращается (см. пояснения ниже).  
		\item Для $x_2$: коэффициенты:
		\[
		\begin{array}{ll}
			\text{Ограничение 1:} & 0,\\[1mm]
			\text{Ограничение 2:} & 1,\\[1mm]
			\text{Ограничение 3:} & 7,\\[1mm]
			\text{Ограничение 4:} & 3.
		\end{array}
		\]
		Переменная $x_2$ имеет ограничение $x_2\ge 0$, поэтому:
		\[
		0\cdot y_1 + 1\cdot y_2 + 7y_3 + 3y_4 \le 8,
		\]
		так как $c_2=8$.
		\item Для $x_3$: коэффициенты:
		\[
		\begin{array}{ll}
			\text{Ограничение 1:} & 4,\\[1mm]
			\text{Ограничение 2:} & 7,\\[1mm]
			\text{Ограничение 3:} & -3,\\[1mm]
			\text{Ограничение 4:} & 8.
		\end{array}
		\]
		Переменная $x_3$ не имеет ограничений по знаку (свободна), поэтому требуемое условие имеет вид:
		\[
		4y_1 + 7y_2 - 3y_3 + 8y_4 = 10,
		\]
		поскольку $c_3=10$.
	\end{itemize}
	
	Таким образом, двойственная задача примет вид:
	\[
	\begin{array}{rl}
		\text{maximize} & y_1 - 2y_3 + 3y_4\\[1mm]
		\text{s.t.} & -4y_1 + 5y_2 + 8y_3 + 6y_4 \le -9,\\[1mm]
		& y_2 + 7y_3 + 3y_4 \le 8,\\[1mm]
		& 4y_1 + 7y_2 - 3y_3 + 8y_4 = 10,\\[1mm]
		& y_1,\, y_3,\, y_4 \ge 0,\quad y_2\in\mathbb{R}.
	\end{array}
	\]
	\vspace{2mm}
	
	\textbf{Замечание.} Полученная двойственная задача составлена на основании стандартных соотношений между переменными исходной задачи и двойственными переменными. Детали перехода можно получить через форму записи лагранжиана для задачи минимизации.
	
	\bigskip
	
	\subsection*{Задача 3}
	
	\medskip
	
	\textbf{} Система
	\[
	A_1x\le b_1,\quad A_2x = b_2
	\]
	совместна тогда и только тогда, если двойственная задача (полученная после введения искусственной целевой функции) совместна. При этом, согласно теореме о двойственности, система не совместна, если существует набор двойственных множителей $\lambda_1,\lambda_2$ (при этом $\lambda_1\ge 0$, $\lambda_2$ свободны), удовлетворяющий:
	\[
	\lambda_1^T A_1 + \lambda_2^T A_2 = 0 \quad \text{и} \quad \lambda_1^T b_1 + \lambda_2^T b_2 < 0.
	\]
	Таким образом, наличие такого двойственного сертификата (разделяющей гиперплоскости) является необходимым и достаточным условием несовместности исходной системы.
	
	\bigskip
	
	\subsection*{Задача 4 }
	
	Введём двойственные переменные:
	\begin{itemize}[leftmargin=1cm]
		\item $u$ для первого (неравенственного) ограничения, $u\ge0$;
		\item $v$ для равенства, $v\in\mathbb{R}$.
	\end{itemize}
	
	Запишем лагранжиан:
	\[
	L(x,u,v)=c_1x_1+c_2x_2 + u\Bigl(b_1-A_{11}x_1-A_{12}x_2\Bigr) + v\Bigl(b_2-A_{21}x_1-A_{22}x_2\Bigr).
	\]
	Перегруппируем по переменным:
	\[
	L(x,u,v)=u\,b_1+v\,b_2 + x_1\Bigl(c_1-A_{11}u-A_{21}v\Bigr)+ x_2\Bigl(c_2-A_{12}u-A_{22}v\Bigr).
	\]
	Чтобы супремум по $x$ был конечен, необходимо:
	\begin{itemize}[leftmargin=1cm]
		\item Для $x_1$, имеющей ограничение $x_1\le 0$, условие
		\[
		c_1-A_{11}u-A_{21}v \ge 0,
		\]
		поскольку при $x_1\le 0$ супремум достигается, если множитель перед $x_1$ неотрицательный.
		\item Для свободной переменной $x_2$:
		\[
		c_2-A_{12}u-A_{22}v=0.
		\]
	\end{itemize}
	Таким образом, двойственная задача имеет вид:
	\[
	\begin{array}{rl}
		\text{minimize} & u\,b_1+v\,b_2\\[1mm]
		\text{s.t.} & A_{11}u+A_{21}v \le c_1,\\[1mm]
		& A_{12}u+A_{22}v = c_2,\\[1mm]
		& u\ge 0,\quad v\in\mathbb{R}.
	\end{array}
	\]
	
	\medskip
	
	Пусть $(x_1,x_2)$ удовлетворяет ограничениям исходной задачи, а $(u,v)$ удовлетворяет ограничениям двойственной задачи, т.е.:
	\[
	\begin{array}{rcl}
		A_{11}x_1+A_{12}x_2 &\le& b_1,\quad A_{21}x_1+A_{22}x_2 = b_2,\quad x_1\le 0,\\[1mm]
		A_{11}u+A_{21}v &\le& c_1,\quad A_{12}u+A_{22}v = c_2,\quad u\ge0.
	\end{array}
	\]
	Умножая первое ограничение исходной задачи на $u$ (с $u\ge0$) и второе на $v$, затем суммируя, получаем:
	\[
	u(A_{11}x_1+A_{12}x_2) + v(A_{21}x_1+A_{22}x_2) \le u\,b_1 + v\,b_2.
	\]
	С другой стороны, с учетом равенств в двойственной задаче можно показать, что
	\[
	c_1x_1+c_2x_2 \le u\,b_1+v\,b_2.
	\]
	Это и есть неравенство слабой двойственности.
	
	\medskip
	
	Из теоремы сильной двойственности известно, что если существуют допустимые решения $(x_1,x_2)$ и $(u,v)$, для которых выполняется равенство
	\[
	c_1x_1+c_2x_2 = u\,b_1+v\,b_2,
	\]
	то эти решения оптимальны для соответствующих задач. Обратно, если $(x_1,x_2)$ --- оптимальное решение исходной задачи, а $(u,v)$ --- оптимальное решение двойственной, то значения целевых функций равны. Таким образом, условие равенства функционалов является необходимым и достаточным условием оптимальности.
	
	\newpage
	
	\section*{Дополнительная часть}
	
	\subsection*{Задача 1}
	
	
	\textbf{}
	
	Введём двойственные переменные $\lambda_i\in\mathbb{R}$ для ограничений
	\[
	a_i^T x + y_i = b_i,\quad i=1,\dots,m.
	\]
	Запишем лагранжиан:
	\[
	L(x,y,\lambda)= -\sum_{i=1}^{m}\ln y_i + \sum_{i=1}^{m}\lambda_i\Bigl(b_i-a_i^T x-y_i\Bigr).
	\]
	Перегруппируем:
	\[
	L(x,y,\lambda)= \Bigl(\sum_{i=1}^{m}\lambda_i b_i\Bigr) - \Bigl(\sum_{i=1}^{m}\lambda_i a_i^T\Bigr)x + \sum_{i=1}^{m}\Bigl[-\ln y_i-\lambda_i y_i\Bigr].
	\]
	Чтобы функционал был ограничен сверху по $x$, необходимо
	\[
	\sum_{i=1}^{m}\lambda_i a_i = 0.
	\]
	Далее, рассмотрим оптимизацию по $y_i>0$. Для каждого $i$ решаем задачу
	\[
	\sup_{y_i>0}\Bigl[-\ln y_i-\lambda_i y_i\Bigr].
	\]
	Оптимальное значение находится по условию первого порядка. Дифференцируя по $y_i$:
	\[
	-\frac{1}{y_i}-\lambda_i=0 \quad\Longrightarrow\quad y_i=-\frac{1}{\lambda_i}.
	\]
	Чтобы оптимальное $y_i$ оказалось положительным, должно выполняться условие $\lambda_i<0$. Подставляя найденное $y_i$, получаем:
	\[
	-\ln\Bigl(-\frac{1}{\lambda_i}\Bigr)-\lambda_i\Bigl(-\frac{1}{\lambda_i}\Bigr)
	= -\ln\Bigl(-\frac{1}{\lambda_i}\Bigr)+1 = 1+\ln(-\lambda_i).
	\]
	Таким образом, двойственная функция равна (с учётом суммирования по $i$):
	\[
	g(\lambda)= \sum_{i=1}^{m}\lambda_i b_i + \sum_{i=1}^{m}\Bigl(1+\ln(-\lambda_i)\Bigr),
	\]
	при условии
	\[
	\sum_{i=1}^{m}\lambda_i a_i =0,\quad \lambda_i<0,\quad i=1,\dots,m.
	\]
	Отметим, что постоянный член $\sum_{i=1}^{m}1=m$ можно опустить при формулировке оптимизации, если добавить его в значение целевой функции.
	
	\medskip
	
	Таким образом, двойственная задача имеет вид:
	\[
	\begin{array}{rl}
		\text{maximize} & \displaystyle \sum_{i=1}^{m}\lambda_i b_i + \sum_{i=1}^{m}\ln(-\lambda_i) + m\\[2mm]
		\text{s.t.} & \displaystyle \sum_{i=1}^{m}\lambda_i a_i =0,\\[1mm]
		& \lambda_i<0,\quad i=1,\dots,m.
	\end{array}
	\]
	
	\bigskip
	
	\subsection*{Задача 2}

	
	\medskip
	
	\textbf{LP-релаксация:}
	\[
	\begin{array}{rl}
		\text{minimize} & c^T x\\[1mm]
		\text{s.t.} & Ax \le b,\\[1mm]
		& 0 \le x_i \le 1,\quad i=1,\dots,d.
	\end{array}
	\]
	Пусть в этой задаче зафиксированы дополнительные ограничения на $x$; стандартное построение двойственной задачи для LP с ограничениями типа $Ax\le b$ и коробочными ограничениями ($0\le x_i\le 1$) приводит к системе двойственных ограничений, включающей дополнительные двойственные переменные для ограничений $x_i\le1$. Обобщённо, двойственная задача может быть записана в виде:
	\[
	\begin{array}{rl}
		\text{maximize} & b^T u + \mathbf{1}^T v\\[1mm]
		\text{s.t.} & A^T u + w - v = c,\\[1mm]
		& u\ge 0,\quad v\ge 0,\quad w\ge 0,
	\end{array}
	\]
	где $u$ соответствует ограничениям $Ax\le b$, а $v$ и $w$ --- двойственным переменным для ограничений $x_i\le 1$ и $-x_i\le 0$ соответственно.
	
	\medskip
	
	\textbf{Вторая релаксация.}  
	Запишем задачу в эквивалентном виде:
	\[
	\begin{array}{rl}
		\text{minimize} & c^T x\\[1mm]
		\text{s.t.} & Ax \le b,\\[1mm]
		& x_i(1-x_i)=0,\quad i=1,\dots,d.
	\end{array}
	\]
	В данной записи нелинейные условия $x_i(1-x_i)=0$ гарантируют, что $x_i\in\{0,1\}$. Для построения двойственной задачи применяют лагранжево–Штурмову релаксацию: для каждого условия вводятся мультипликаторы $\lambda_i$ (соответствующей знаковой характеристикой). Тогда соответствующая двойственная задача (называемая Лагранжевой релаксацией) имеет вид
	\[
	\begin{array}{rl}
		\text{maximize} & \tilde{d}(u,\lambda)\\[2mm]
		\text{s.t.} & \text{условия на двойственные переменные } u,\, \lambda,
	\end{array}
	\]
	где нижняя оценка на оптимальное значение исходной задачи получается за счёт оптимизации по $(u,\lambda)$.
	

	
	\bigskip
	

\end{document}
